\chapter{API REST}
\label{chap:api}

\section{Principi di design}
\label{sec:api_design}

Le API REST del backend seguono le seguenti convenzioni:

\begin{itemize}[leftmargin=*]
    \item \textbf{Autenticazione}: JWT Bearer token HS256 con scadenza di 30~minuti,
          passato nell'header \texttt{Authorization: Bearer \{token\}}.
    \item \textbf{Formato degli errori}: struttura uniforme
          \texttt{\{"error": "...", "code": "...", "details": \{...\}\}}.
    \item \textbf{Validazione}: Pydantic v2 valida automaticamente tutti gli schema;
          la violazione produce un \texttt{400 Bad Request} con dettaglio dei campi.
    \item \textbf{Paginazione}: tutti gli endpoint di listing espongono i parametri
          \texttt{limit} (1--100, default 20) e \texttt{offset} ($\geq 0$, default 0).
    \item \textbf{Documentazione}: Swagger UI automatica all'endpoint \texttt{/docs}.
\end{itemize}

\section{Autenticazione (\texttt{/auth})}
\label{sec:api_auth}

\begin{table}[ht]
\centering
\small
\begin{tabularx}{\linewidth}{l l l X}
\toprule
\textbf{Metodo} & \textbf{Path} & \textbf{Auth} & \textbf{Descrizione} \\
\midrule
POST & \texttt{/auth/register} & No  & Registrazione. Il primo utente del dominio diventa \texttt{DOMAIN\_ADMIN}. \\
POST & \texttt{/auth/login}    & No  & Restituisce JWT firmato HS256. \\
GET  & \texttt{/auth/me}       & JWT & Profilo dell'utente autenticato. \\
POST & \texttt{/auth/logout}   & JWT & Invalidazione lato client (stateless). \\
\bottomrule
\end{tabularx}
\caption{Endpoint di autenticazione.}
\end{table}

\section{Cartelle (\texttt{/folders})}
\label{sec:api_folders}

\begin{table}[ht]
\centering
\small
\begin{tabularx}{\linewidth}{l l l X}
\toprule
\textbf{Metodo} & \textbf{Path} & \textbf{Auth} & \textbf{Descrizione} \\
\midrule
POST  & \texttt{/folders}             & JWT & Crea cartella (con \texttt{parent\_id} opzionale). \\
GET   & \texttt{/folders}             & JWT & Lista cartelle root o figlie di \texttt{parent\_id}. \\
GET   & \texttt{/folders/\{id\}}      & JWT & Dettagli cartella. \\
DELETE & \texttt{/folders/\{id\}}     & JWT & Cancellazione in cascata (DB + S3). \\
PATCH & \texttt{/folders/\{id\}/move} & JWT & Spostamento a nuovo \texttt{parent\_id}. \\
POST  & \texttt{/folders/\{id\}/copy} & JWT & Copia ricorsiva (struttura + file S3). \\
\bottomrule
\end{tabularx}
\caption{Endpoint per la gestione delle cartelle.}
\end{table}

\section{Documenti (\texttt{/documents})}
\label{sec:api_documents}

\begin{table}[ht]
\centering
\small
\begin{tabularx}{\linewidth}{l l l X}
\toprule
\textbf{Metodo} & \textbf{Path} & \textbf{Auth} & \textbf{Descrizione} \\
\midrule
POST & \texttt{/documents/upload}         & JWT & Crea record + restituisce presigned URL S3. \\
POST & \texttt{/documents/upload-ai}      & JWT & Upload asincrono con classificazione AI (202 Accepted). \\
POST & \texttt{/documents/\{id\}/confirm} & JWT & Finalizza upload: imposta \texttt{size\_bytes}. \\
GET  & \texttt{/documents}                & JWT & Lista documenti (con paginazione). \\
GET  & \texttt{/documents/\{id\}}         & JWT & Metadati documento. \\
GET  & \texttt{/documents/\{id\}/download}& JWT & Download in streaming (\texttt{StreamingResponse}). \\
DELETE & \texttt{/documents/\{id\}}       & JWT & Cancellazione DB + S3. \\
\bottomrule
\end{tabularx}
\caption{Endpoint per la gestione dei documenti.}
\end{table}

\subsection{Pattern two-phase upload}

Il flusso di caricamento diretto si articola in tre passi:
\begin{enumerate}[leftmargin=*]
    \item \texttt{POST /documents/upload}: il backend crea un \texttt{Document} con
          \texttt{size\_bytes = NULL} e genera un presigned URL S3 con scadenza 1~ora.
    \item Il client esegue un HTTP PUT direttamente su SeaweedFS (il dato non passa dal backend).
    \item \texttt{POST /documents/\{id\}/confirm} con \texttt{\{"size\_bytes": N\}}:
          il backend finalizza il record rendendolo visibile.
\end{enumerate}

\section{Permessi (\texttt{/permissions})}
\label{sec:api_permissions}

\begin{table}[ht]
\centering
\small
\begin{tabularx}{\linewidth}{l l l X}
\toprule
\textbf{Metodo} & \textbf{Path} & \textbf{Auth} & \textbf{Descrizione} \\
\midrule
POST   & \texttt{/permissions}        & JWT & Condividi risorsa per \texttt{user\_id}. \\
POST   & \texttt{/permissions/share}  & JWT & Condividi risorsa per email (VIEWER o EDITOR). \\
GET    & \texttt{/permissions}        & JWT & Lista permessi concessi all'utente corrente. \\
DELETE & \texttt{/permissions/\{id\}} & JWT & Revoca permesso (solo il proprietario). \\
\bottomrule
\end{tabularx}
\caption{Endpoint per la gestione dei permessi.}
\end{table}

La condivisione \`e consentita solo tra utenti dello stesso dominio email. Il livello
\texttt{VIEWER} garantisce accesso in sola lettura; \texttt{EDITOR} permette anche la
scrittura (upload nella cartella, rinomina).

\section{Account email (\texttt{/email-accounts})}
\label{sec:api_email}

\begin{table}[ht]
\centering
\small
\begin{tabularx}{\linewidth}{l l l X}
\toprule
\textbf{Metodo} & \textbf{Path} & \textbf{Auth} & \textbf{Descrizione} \\
\midrule
POST   & \texttt{/email-accounts}              & JWT & Aggiunge account IMAP (credenziali cifrate). \\
GET    & \texttt{/email-accounts}              & JWT & Lista account dell'utente. \\
GET    & \texttt{/email-accounts/\{id\}}       & JWT & Dettagli account. \\
PATCH  & \texttt{/email-accounts/\{id\}}       & JWT & Aggiorna configurazione. \\
DELETE & \texttt{/email-accounts/\{id\}}       & JWT & Cancella account e job associati. \\
POST   & \texttt{/email-accounts/\{id\}/test}  & JWT & Testa connessione IMAP. \\
\bottomrule
\end{tabularx}
\caption{Endpoint per la gestione degli account email.}
\end{table}

Le credenziali (password o token OAuth2) sono cifrate con Fernet prima del salvataggio
e non compaiono mai nelle risposte API.

\section{Layer agentico (\texttt{/agent})}
\label{sec:api_agent}

\begin{table}[ht]
\centering
\small
\begin{tabularx}{\linewidth}{l l l X}
\toprule
\textbf{Metodo} & \textbf{Path} & \textbf{Auth} & \textbf{Descrizione} \\
\midrule
GET  & \texttt{/agent/operations}                    & JWT & Timeline operazioni agentiche. \\
GET  & \texttt{/agent/proposals}                     & JWT & Inbox delle proposte in attesa di revisione. \\
POST & \texttt{/agent/proposals/\{id\}/accept}       & JWT & Accetta la proposta dell'agente. \\
POST & \texttt{/agent/proposals/\{id\}/move}         & JWT & Sposta nella cartella scelta dall'utente. \\
POST & \texttt{/agent/proposals/\{id\}/reject}       & JWT & Rigetta la proposta. \\
\bottomrule
\end{tabularx}
\caption{Endpoint per il layer agentico.}
\end{table}

\section{Ricerca (\texttt{/search})}
\label{sec:api_search}

\begin{lstlisting}[language=bash, caption={Esempio di query di ricerca.}]
GET /search?q=fattura+fornitore+2024&limit=20&offset=0
Authorization: Bearer {token}
\end{lstlisting}

La ricerca combina full-text search (campo \texttt{text} con analizzatore
\texttt{standard}) con ricerca vettoriale k-NN sul campo \texttt{embedding}.
I risultati sono ordinati per rilevanza semantica.

\section{Admin (\texttt{/admin})}
\label{sec:api_admin}

\begin{table}[ht]
\centering
\small
\begin{tabularx}{\linewidth}{l l l X}
\toprule
\textbf{Metodo} & \textbf{Path} & \textbf{Auth} & \textbf{Descrizione} \\
\midrule
GET   & \texttt{/admin/metrics}              & JWT + DOMAIN\_ADMIN & Statistiche aggregate del dominio. \\
PATCH & \texttt{/admin/users/\{id\}/status}  & JWT + DOMAIN\_ADMIN & Abilita/disabilita utente. \\
PATCH & \texttt{/admin/users/\{id\}/role}    & JWT + DOMAIN\_ADMIN & Cambia ruolo utente. \\
\bottomrule
\end{tabularx}
\caption{Endpoint di amministrazione del dominio.}
\end{table}

\section{Health check (\texttt{/health})}
\label{sec:api_health}

\begin{lstlisting}[language=Python, caption={Implementazione dell'endpoint \texttt{/health}.}]
@router.get("/health")
def health(db: Session = Depends(get_db)):
    checks = {"status": "healthy", "database": "unknown", "s3": "unknown"}
    try:
        db.execute(text("SELECT 1"))
        checks["database"] = "ok"
    except Exception as e:
        checks["database"] = "error"
        checks["status"] = "unhealthy"
    try:
        s3 = get_s3_client()
        checks["s3"] = "ok"
    except Exception as e:
        checks["s3"] = "error"
        # S3 degraded non rende l'intero servizio unhealthy
    status_code = 200 if checks["status"] == "healthy" else 503
    return JSONResponse(checks, status_code=status_code)
\end{lstlisting}

L'endpoint distingue un errore critico (database irraggiungibile $\rightarrow$ 503) da
uno non critico (S3 degradato $\rightarrow$ warning ma 200), in coerenza con l'impatto
reale sul servizio: il read path dei metadati rimane funzionante anche senza S3.
